% TV1--TV3 cung cấp lệnh/log; TV4 xác nhận số liệu; TV5 tổng hợp.
\chapter{Khả năng tái lập và đối chiếu yêu cầu}
\label{app:reproducibility}
\section{Nguồn, phiên bản và lệnh chạy}
\todo{Điền URL nguồn, license/nguồn hình nếu cần, commit tác giả, commit nhóm, environment lock, cấu hình và lệnh thực tế đã chạy. Không để lệnh dự kiến như thể đã kiểm chứng.}
\begin{table}[htbp]\centering\small
\caption{Bằng chứng đủ để người khác chạy lại.}\label{tab:reproducibility}
\begin{tabularx}{\linewidth}{@{}P{42mm}Y@{}}\toprule
Đường dẫn README Source & \blank\\
Lệnh tạo môi trường & \blank\\
Lệnh smoke test & \blank\\
Lệnh chạy panel chính & \blank\\
Lệnh dựng bảng / hình & \blank\\
CSV / config / log / checkpoint & \blank\\
Commit số liệu PDF--video & \blank\\
Kết quả thử môi trường sạch & \blank\\
Lệnh render Manim / assets & \blank\\
Video / phụ đề / url.txt & \blank\\\bottomrule
\end{tabularx}\end{table}

\section{Giao diện kết quả tối thiểu}
\noindent Mỗi dòng CSV là một cấu hình, checkpoint và seed. Mỗi dòng có \codepath{run_id} duy nhất. MLP để trống \codepath{G,k} và lưu \codepath{config_id}, \codepath{activation} để phân biệt các cấu hình cùng kiến trúc. Schema chuẩn 22 cột trong \codepath{results/runs.csv} (đồng bộ với \codepath{results_schema.md}):
\begin{quote}\small\ttfamily\raggedright
run\_id, experiment\_id, task, model, config\_id, activation, architecture, G, k, seed, split\_id, params, train\_rmse, val\_rmse, test\_rmse, time\_s, optimizer, steps, max\_iter, closure\_evals, config\_path, code\_commit
\end{quote}
\todo{Đính kèm cấu hình phần cứng/phiên bản chạy thực tế sau khi ghép nhóm; ghi quy tắc ánh xạ các \codepath{run_id} vào một điểm trung bình và độ lệch chuẩn trong tệp \codepath{results/summary.json}.}

\section{Đối chiếu yêu cầu với vị trí báo cáo}
\begin{longtable}{@{}P{64mm}P{51mm}P{18mm}@{}}
\caption{Bảng kiểm cuối kỳ; điền trạng thái sau khi kiểm chứng.}\label{tab:traceability}\\
\toprule\textbf{Yêu cầu} & \textbf{Vị trí dự kiến} & \textbf{Đã kiểm}\\\midrule\endfirsthead
\toprule\textbf{Yêu cầu} & \textbf{Vị trí dự kiến} & \textbf{Đã kiểm}\\\midrule\endhead
Bài toán, ví dụ, hai hướng trước, đóng góp phân loại & Chương~\ref{ch:introduction} & \blank\\
Ký hiệu, sơ đồ, cơ chế, định lý và giả định & Chương~\ref{ch:theory} & \blank\\
Huấn luyện, lưới, độ phức tạp, liên hệ môn học & Chương~\ref{ch:training} & \blank\\
Ánh xạ công thức--code và tác động thay đổi & Bảng~\ref{tab:code-mapping} & \blank\\
Thực nghiệm tác giả, setup, dữ liệu, thay đổi & Chương~\ref{ch:experiments} & \blank\\
Ít nhất 4 điểm KAN và 4 MLP; cùng split; run ID & Bảng~\ref{tab:core-results}, Hình~\ref{fig:rmse-params} & \blank\\
Khảo sát G, train/test, tham số và thời gian & Hình~\ref{fig:rmse-grid}, \ref{fig:rmse-time} & \blank\\
Đối chiếu đúng panel, phân tích sai khác & Bảng~\ref{tab:paper-comparison} & \blank\\
Hai điểm mạnh, hai hạn chế, cải tiến có căn cứ & Chương~\ref{ch:evaluation} & \blank\\
Đóng góp, tỷ lệ thực tế, AI, tài liệu tham khảo & Phụ lục~\ref{app:contributions} & \blank\\
README/lệnh/commit, thử môi trường sạch & Phụ lục~\ref{app:reproducibility} & \blank\\
Video 5--10 phút, Manim, phụ đề, nguồn, công khai & URL video và nguồn Manim & \blank\\
Source, Manim, Report.pdf, url.txt và tên MSSV & Gói nộp cuối cùng & \blank\\\bottomrule
\end{longtable}
\writingnote{Video dùng hình tự dựng, nguồn hiện lúc sử dụng, phụ đề xuyên suốt; đối chiếu số liệu với PDF và CSV. Kế hoạch yêu cầu hashtag \#fithcmus \#machinelearning \#ai \#ml. Kiểm yêu cầu nộp theo đề gốc trước khi đóng gói.}
